\newpage
\section*{Part C: Solved Examples 6 to 10 (Section 7.3, Section Formula)}

%---------------------------------------------------------------
\subsection*{Example 6}
\textbf{Question.} Find the coordinates of the point which divides the line segment joining the points $(4,-3)$ and $(8,5)$ in the ratio $3:1$ internally.

\textbf{Solution.} Let $A(4,-3)$, $B(8,5)$ and let $P(x,y)$ divide $AB$ internally in the ratio $m_1:m_2=3:1$. So $(x_1,y_1)=(4,-3)$, $(x_2,y_2)=(8,5)$, $m_1=3$, $m_2=1$.

By the section formula, $P\left(\dfrac{m_1x_2+m_2x_1}{m_1+m_2},\dfrac{m_1y_2+m_2y_1}{m_1+m_2}\right)$:
\begin{align*}
x&=\frac{3(8)+1(4)}{3+1}=\frac{24+4}{4}=\frac{28}{4}=7\\
y&=\frac{3(5)+1(-3)}{3+1}=\frac{15-3}{4}=\frac{12}{4}=3
\end{align*}
\slip{$m_1$ multiplies the $x_2$ (and $y_2$) of the \emph{far} point $B$, and $m_2$ multiplies the near point $A$. Swapping them gives the point that divides in the ratio $1:3$ instead.}

\ans{The required point is $(7,3)$.}

%---------------------------------------------------------------
\subsection*{Example 7}
\textbf{Question.} In what ratio does the point $(-4,6)$ divide the line segment joining the points $A(-6,10)$ and $B(3,-8)$?

\textbf{Solution.} Let $P(-4,6)$ divide $AB$ internally in the ratio $m_1:m_2$. Here $(x_1,y_1)=(-6,10)$ and $(x_2,y_2)=(3,-8)$.

By the section formula $P\left(\dfrac{m_1x_2+m_2x_1}{m_1+m_2},\dfrac{m_1y_2+m_2y_1}{m_1+m_2}\right)$:
\[
(-4,6)=\left(\frac{m_1(3)+m_2(-6)}{m_1+m_2},\ \frac{m_1(-8)+m_2(10)}{m_1+m_2}\right)=\left(\frac{3m_1-6m_2}{m_1+m_2},\ \frac{-8m_1+10m_2}{m_1+m_2}\right).
\]
Equating the $x$-coordinates:
\begin{align*}
\frac{3m_1-6m_2}{m_1+m_2}&=-4\\
3m_1-6m_2&=-4(m_1+m_2)=-4m_1-4m_2\\
3m_1+4m_1&=6m_2-4m_2\\
7m_1&=2m_2\\
\frac{m_1}{m_2}&=\frac27
\end{align*}
So the ratio is $m_1:m_2=2:7$.

\textbf{Check with the $y$-coordinates.} Put $m_1=2$, $m_2=7$:
\[\frac{-8(2)+10(7)}{2+7}=\frac{-16+70}{9}=\frac{54}{9}=6,\]
which is the given $y$-coordinate. So the ratio is correct.

\ans{The point $(-4,6)$ divides $AB$ internally in the ratio $2:7$.}

%---------------------------------------------------------------
\subsection*{Example 8}
\textbf{Question.} Find the coordinates of the points of trisection (i.e., points dividing in three equal parts) of the line segment joining the points $A(2,-2)$ and $B(-7,4)$.

\textbf{Solution.} Let $P$ and $Q$ be the points of trisection, so $AP=PQ=QB$ with $P$ nearer to $A$. Here $(x_1,y_1)=(2,-2)$ and $(x_2,y_2)=(-7,4)$.

$P$ is one part from $A$ and two parts from $B$, so it divides $AB$ in the ratio $1:2$. $Q$ is two parts from $A$ and one part from $B$, so it divides $AB$ in the ratio $2:1$.
By the section formula $\left(\dfrac{m_1x_2+m_2x_1}{m_1+m_2},\dfrac{m_1y_2+m_2y_1}{m_1+m_2}\right)$:

For $P$ with $m_1:m_2=1:2$:
\[
x=\frac{1(-7)+2(2)}{1+2}=\frac{-7+4}{3}=\frac{-3}{3}=-1,\qquad y=\frac{1(4)+2(-2)}{1+2}=\frac{4-4}{3}=\frac{0}{3}=0 .
\]
For $Q$ with $m_1:m_2=2:1$:
\[
x=\frac{2(-7)+1(2)}{2+1}=\frac{-14+2}{3}=\frac{-12}{3}=-4,\qquad y=\frac{2(4)+1(-2)}{2+1}=\frac{8-2}{3}=\frac{6}{3}=2 .
\]

\textbf{Verification.} By the mid-point formula $\left(\dfrac{x_1+x_2}{2},\dfrac{y_1+y_2}{2}\right)$, the mid-point of $PB$ is
\[
\left(\dfrac{-1+(-7)}{2},\dfrac{0+4}{2}\right)=(-4,2)=Q,
\]
as it should be because $Q$ is halfway between $P$ and $B$.

\ans{The points of trisection are $P(-1,0)$ and $Q(-4,2)$.}

%---------------------------------------------------------------
\subsection*{Example 9}
\textbf{Question.} Find the ratio in which the $y$-axis divides the line segment joining the points $(5,-6)$ and $(-1,-4)$. Also find the point of intersection.

\textbf{Solution.} Let $A(5,-6)$, $B(-1,-4)$, and let the $y$-axis divide $AB$ in the ratio $k:1$ at the point $P$. So $(x_1,y_1)=(5,-6)$, $(x_2,y_2)=(-1,-4)$, $m_1=k$, $m_2=1$.

By the section formula,
\[
P=\left(\frac{k(-1)+1(5)}{k+1},\ \frac{k(-4)+1(-6)}{k+1}\right)=\left(\frac{5-k}{k+1},\ \frac{-4k-6}{k+1}\right).
\]
A point on the $y$-axis has $x$-coordinate $0$:
\[
\frac{5-k}{k+1}=0\ \Rightarrow\ 5-k=0\ \Rightarrow\ k=5 .
\]
So the ratio is $5:1$. Put $k=5$ in the $y$-coordinate:
\[y=\frac{-4(5)-6}{5+1}=\frac{-20-6}{6}=\frac{-26}{6}=-\frac{13}{3}.\]
\slip{On the $y$-axis it is the $x$-coordinate that is $0$. Setting the $y$-coordinate to $0$ here would give a different, wrong ratio.}

\ans{The $y$-axis divides the segment in the ratio $5:1$, and the point of intersection is $\left(0,-\dfrac{13}{3}\right)$.}

%---------------------------------------------------------------
\subsection*{Example 10}
\textbf{Question.} If the points $A(6,1)$, $B(8,2)$, $C(9,4)$ and $D(p,3)$ are the vertices of a parallelogram, taken in order, find the value of $p$.

\textbf{Solution.} The diagonals of a parallelogram bisect each other, so the mid-point of $AC$ equals the mid-point of $BD$.

By the mid-point formula $\left(\dfrac{x_1+x_2}{2},\dfrac{y_1+y_2}{2}\right)$:
\[
\begin{gathered}
\text{Mid-point of }AC=\left(\frac{6+9}{2},\frac{1+4}{2}\right)=\left(\frac{15}{2},\frac52\right),\\
\text{Mid-point of }BD=\left(\frac{8+p}{2},\frac{2+3}{2}\right)=\left(\frac{8+p}{2},\frac52\right).
\end{gathered}
\]
Equating the $x$-coordinates: $\dfrac{15}{2}=\dfrac{8+p}{2}$, so $8+p=15$ and $p=7$.
The $y$-coordinates are both $\dfrac52$, which is consistent.

\textbf{Verification.} With $D(7,3)$, opposite sides should be equal. By the distance formula $d^2=(x_2-x_1)^2+(y_2-y_1)^2$:
\begin{align*}
AB^2&=(8-6)^2+(2-1)^2=2^2+1^2=4+1=5\\
DC^2&=(9-7)^2+(4-3)^2=2^2+1^2=4+1=5\\
BC^2&=(9-8)^2+(4-2)^2=1^2+2^2=1+4=5\\
AD^2&=(7-6)^2+(3-1)^2=1^2+2^2=1+4=5
\end{align*}
So $AB=DC$ and $BC=AD$: opposite sides are equal, as in a parallelogram.

\ans{$p=7$.}
